% !TEX root = ../main.tex
\section{Proof of Theorem~\ref{thm:upper-beta0}}
\label{app:upper-beta0}

All references to~\cite{li2026smoothing} in this appendix use the numbering
of arXiv:2602.14185v2. Let
\[
\ell:=\max\Bigl\{L_{\calD,\bx},G_{\calA},\frac1\rho\Bigr\},
\qquad r:=\frac1\gamma,\qquad \delta:=\frac1s,
\qquad \mathcal Y:=\{\blmbd\in\RR^m:\|\blmbd\|\le2L_h\},
\qquad D_{\mathcal Y}:=4L_h,
\]
where \(L_{\calD,\bx}\) is defined in~\eqref{eq:x-weak-convexity-modulus}.
We take \(L_h>0\). If the supplied class constant is zero, we use the
admissible Lipschitz bound \(L_h=1/16\) throughout this appendix; the
choice \(c_s=\max\{1,16L_h\}\) below remains \(1\), and the algorithm is
unchanged. All auxiliary constants are chosen from the fixed parameters
and class bounds, independently of the instance and its dimensions.

\paragraph{Step 1: reduction to Perturbed Smoothed GDA}
Since \(\beta=0\) and \(\bmu^0=\zero\), we have \(\bmu^k=\zero\) for all
\(k\). The potential function becomes
\[
\mathcal P_s(\bx,\blmbd;\bz,\zero)
=\calD_\rho(\bx,\blmbd)+\frac r2\|\bx-\bz\|^2
-\frac\delta2\|\blmbd\|^2.
\]
The correspondence with~\cite[(3.1) and Algorithm~3]{li2026smoothing} is
\[
f_{\mathrm{Li}}=\calD_\rho,\quad
(r_x,r_y^t,c^t,\alpha^t_{\mathrm{Li}},\beta^t_{\mathrm{Li}})
=(r,\delta,\eta,\tau,\alpha),\quad
\mathcal X=\RR^n.
\]
In particular, both algorithms use the old multiplier in the primal step,
the new primal iterate in the dual step, and the new primal iterate in the
center step. By~\eqref{eq:prox-multiplier}, the dual update reads
\[
\blmbd^{k+1}
=\Bigl(1-\frac\tau\rho-\frac\tau s\Bigr)\blmbd^k
+\frac\tau\rho\bq^k,
\qquad \|\bq^k\|\le L_h.
\]
The step-size conditions in Step~3 give
\(\tau\le1/(11(\ell+1/s))<1/(1/\rho+1/s)\).
Thus both coefficients are nonnegative and sum to \(1-\tau/s\le1\).
Induction yields \(\|\blmbd^k\|\le L_h\). Every multiplier lies in the
interior of \(\mathcal Y\), so the projection onto \(\mathcal Y\) does not
change the iterates and its normal cone at \(\blmbd^k\) is
\(\{\zero\}\). The primal projection is the identity. Consequently,
SPLALM generates exactly the sequence of Perturbed Smoothed GDA applied to
\(\calD_\rho\) on \(\RR^n\times\mathcal Y\), from zero.

\paragraph{Step 2: assumptions and the unperturbed objective}
The gradient estimates in Step~2 of Appendix~\ref{app:upper-beta1} imply
\[
\begin{aligned}
\|\nabla_{\bx}\calD_\rho(\bx,\blmbd)
-\nabla_{\bx}\calD_\rho(\bx',\blmbd')\|
&\le\ell\bigl(\|\bx-\bx'\|+\|\blmbd-\blmbd'\|\bigr),\\
\|\nabla_{\blmbd}\calD_\rho(\bx,\blmbd)
-\nabla_{\blmbd}\calD_\rho(\bx',\blmbd')\|
&\le\ell\bigl(\|\bx-\bx'\|+\|\blmbd-\blmbd'\|\bigr).
\end{aligned}
\]
Hence Assumption~2.1 of~\cite{li2026smoothing} holds. Moreover,
\(\calD_\rho(\bx,\cdot)\) is concave and \(\mathcal Y\) is a compact
convex set containing zero, as required by its Assumption~2.2.
Taking \(\by=\calA(\bx)\) in~\eqref{eq:reduced-al} gives
\(\calD_\rho(\bx,\blmbd)\le F(\bx)\). For any
\(\blmbd\in\partial h(\calA(\bx))\), proximal optimality gives
\(\by(\bx,\blmbd)=\calA(\bx)\) and equality. Such a multiplier belongs
to \(\mathcal Y\), since its norm is at most \(L_h\). Therefore
\[
\max_{\blmbd\in\mathcal Y}\calD_\rho(\bx,\blmbd)=F(\bx),
\qquad \inf F>-\infty.
\]
The descent argument of~\cite[Section~4.2]{li2026smoothing} uses a lower
bound on the potential function. Accordingly, its minimum objective value
can be replaced by the infimum; attainment of \(\inf F\) is not needed.
The primal subproblems with quadratic regularization attain their minima
by strong convexity and coercivity.

\paragraph{Step 3: step-size conditions}
Set \(a_s:=r-\ell-\delta\) and \(b_s:=\ell+\delta\), and define
\[
\omega_s:=\frac1{\sqrt\delta}\,
\frac{a_s+\tau b_s(3r-2b_s)}{\tau a_s^{3/2}}.
\]
Under the correspondence in Step~1, Condition~3.2(ii)
of~\cite{li2026smoothing} reads
\begin{equation}
\begin{aligned}
&r=\Theta(\ell),\quad r>3\ell,\quad a_s>0,\quad \delta=1/s>0,\\
&\eta=\Theta(1/\ell),\quad 0<\eta<\frac1{r+\ell+\delta},\\
&\tau=\Theta(1/\ell),\quad
0<\tau\le\min\Bigl\{\frac1{11b_s},
\frac{\eta^2a_s^2}{4b_s(1+\eta a_s)^2}\Bigr\},\\
&\alpha=\Theta(\delta/\ell),\quad
0<\alpha\le\min\Bigl\{\frac1{36},
\frac{a_s^2}{384r(r+\ell+\delta)^2},
\frac1{384r\tau\omega_s^2}\Bigr\}.
\end{aligned}
\label{eq:beta0-stepsizes}
\end{equation}
The constants in these order relations are independent of \(s\).
For fixed \(r,\eta,\tau,\ell\) satisfying the first three conditions for
all sufficiently large \(s\), we have \(\omega_s^2=\Theta(s)\).
We can therefore choose a fixed, sufficiently small \(\kappa>0\), as a
function of these fixed quantities, so that
\begin{equation}
\alpha:=\frac\kappa{\ell s}
\label{eq:beta0-alpha}
\end{equation}
satisfies the last condition for all sufficiently large \(s\).
The conditions also give \(\eta<\gamma\) and \(\tau<s\), as required
by Definition~\ref{def:splalm}.

For example, all conditions hold for
\[
\gamma=\frac1{4\ell},\qquad \eta=\frac1{12\ell},\qquad
\tau=\frac1{800\ell},\qquad
\alpha=\frac1{2{,}000{,}000\,\ell s},\qquad
s\ge\max\{1,1/\ell\}.
\]
Indeed, \(2\ell\le a_s\le3\ell\) and \(b_s\le2\ell\), and
\[
\frac{\eta^2a_s^2}{4b_s(1+\eta a_s)^2}
\ge\frac1{392\ell}>\tau,\qquad
\frac{a_s^2}{384r(r+\ell+\delta)^2}\ge\frac1{13824\ell},
\qquad \omega_s^2\le800000\,\ell s.
\]
Thus \(1/(384r\tau\omega_s^2)\ge1/(1536000\,\ell s)>\alpha\),
and the remaining inequalities follow directly.

\paragraph{Step 4: initial gap for the perturbed objective}
Define
\[
\Phi_s(\bx):=\max_{\blmbd\in\mathcal Y}
\Bigl\{\calD_\rho(\bx,\blmbd)-\frac1{2s}\|\blmbd\|^2\Bigr\}.
\]
The maximization in~\eqref{eq:effective-penalty}, with \(\bmu=\zero\),
has an optimal multiplier \(s\bq/(\rho+s)\), where
\(\|\bq\|\le L_h\). This multiplier belongs to \(\mathcal Y\), so
restricting the maximization to \(\mathcal Y\) does not change its value.
It follows that
\begin{equation}
\Phi_s(\bx)=\calD_{\rho+s}(\bx,\zero),\qquad
F(\bx)-\frac{L_h^2}{2(\rho+s)}\le\Phi_s(\bx)\le F(\bx).
\label{eq:beta0-perturbed-objective}
\end{equation}
Let
\[
\begin{aligned}
v_0&:=\calD_\rho(\zero,\zero),\qquad
 d_0:=\min_{\bx}\Bigl\{\calD_\rho(\bx,\zero)
                         +\frac1{2\gamma}\|\bx\|^2\Bigr\},\\
p_0&:=\min_{\bx}\Bigl\{\Phi_s(\bx)+\frac1{2\gamma}\|\bx\|^2\Bigr\}.
\end{aligned}
\]
Since \(\mu_\gamma=1/\gamma-L_{\calD,\bx}>0\), the argument in
Step~4 of Appendix~\ref{app:upper-beta1} gives
\[
d_0\ge v_0-\frac{(G_f+G_{\calA}L_h)^2}{2\mu_\gamma},\qquad
p_0\le F(\zero),\qquad v_0\ge F(\zero)-\frac{L_h^2}{2\rho}.
\]
At zero initialization, the initial gap in
\cite[Definition~3.1(ii)]{li2026smoothing}, with the infimum convention
explained in Step~2, is
\(\Delta_{\Psi_2}=v_0-2d_0+2p_0-\inf\Phi_s\). Hence
\begin{equation}
\begin{aligned}
\Delta_{\Psi_2}
&\le\Delta+\frac{L_h^2}{2\rho}
+\frac{(G_f+G_{\calA}L_h)^2}{\mu_\gamma}
+\frac{L_h^2}{2(\rho+s)}\\
&\le\bar\Delta:=\Delta+\frac{L_h^2}{\rho}
+\frac{(G_f+G_{\calA}L_h)^2}{\mu_\gamma}.
\end{aligned}
\label{eq:beta0-initial-gap}
\end{equation}
The bound \(\bar\Delta\) is uniform over the problem class and independent
of \(s\), the target accuracy, and the dimensions.

\paragraph{Step 5: iteration complexity for game stationarity}
For \(k\ge0\), set
\[
S_{k+1}:=\max\Bigl\{
\|\nabla_{\bx}\calD_\rho(\bx^{k+1},\blmbd^{k+1})\|,
\|\nabla_{\blmbd}\calD_\rho(\bx^{k+1},\blmbd^{k+1})\|\Bigr\}.
\]
Both normal-cone terms vanish by Step~1, so this is the game-stationarity
measure for the unperturbed function \(\calD_\rho\) in
\cite[Theorem~3.3(i)]{li2026smoothing}.
The descent estimate in its Section~4.2 and estimates (4.19)--(4.23),
with~\eqref{eq:beta0-alpha}, give a constant
\(C_*=C_*(\eta\ell,\tau\ell,r/\ell,\kappa)>0\) such that, among the first
\(K\ge1\) updates, some \(0\le k<K\) satisfies
\begin{equation}
S_{k+1}\le C_*\ell\sqrt{\frac{s\bar\Delta}{K}}
+\frac{D_{\mathcal Y}}s.
\label{eq:beta0-gs-rate}
\end{equation}
Here we take \(s\) sufficiently large that \(\ell s\ge1\), so the
term of order \(\sqrt{\ell\bar\Delta/K}\) is absorbed into the first
term. We average the descent estimate over \(k=0,\ldots,K-1\);
a fixed factor in this averaging is included in \(C_*\).
Thus the index in~\eqref{eq:beta0-gs-rate} refers to an output among the
first \(K\) updates of Definition~\ref{def:complexity-measure}.

Choose
\begin{equation}
c_s:=\max\{1,4D_{\mathcal Y}\}=\max\{1,16L_h\},\qquad
s:=\frac{c_s}{\epsilon},\qquad
K:=\max\Bigl\{1,\Bigl\lceil
16C_*^2\ell^2c_s\bar\Delta\,\epsilon^{-3}
\Bigr\rceil\Bigr\}.
\label{eq:beta0-budget}
\end{equation}
Then \(D_{\mathcal Y}/s\le\epsilon/4\), and the first term
in~\eqref{eq:beta0-gs-rate} is at most \(\epsilon/4\).
Consequently, some \(0\le k<K\) has \(S_{k+1}\le\epsilon/2\).

\paragraph{Step 6: from \(S_{k+1}\) to the KKT residual}
Write \(\mathbf d_k:=\blmbd^{k+1}-\blmbd^k\). The dual update and
Step~2 yield
\[
\|\mathbf d_k\|\le\tau\Bigl(S_{k+1}
+\frac{\|\mathbf d_k\|}\rho+\frac{L_h}s\Bigr),\qquad
\|\mathbf d_k\|\le\frac\tau{1-\tau/\rho}
\Bigl(S_{k+1}+\frac{L_h}s\Bigr).
\]
By~\eqref{eq:kkt-residual} and~\eqref{eq:prox-multiplier}, the output
residual is
\[
R_k:=\mathcal R_{\mathrm{KKT}}(\bx^{k+1},\by^{k+1},\bq^k)
=\max\Bigl\{\|\nabla_{\bx}\calD_\rho(\bx^{k+1},\blmbd^k)\|,
\|\nabla_{\blmbd}\calD_\rho(\bx^{k+1},\blmbd^k)\|\Bigr\}.
\]
Both gradients change by at most \(\ell\|\mathbf d_k\|\) when the
multiplier changes from \(\blmbd^{k+1}\) to \(\blmbd^k\). Since
\(\tau\ell\le1/11\) and \(\tau/\rho\le1/11\), we obtain
\begin{equation}
R_k\le\Bigl(1+\frac{\tau\ell}{1-\tau/\rho}\Bigr)S_{k+1}
+\frac{\tau\ell}{1-\tau/\rho}\frac{L_h}s
\le\frac{11}{10}S_{k+1}+\frac{L_h}{10s}.
\label{eq:beta0-residual-transfer}
\end{equation}
For the index selected in Step~5, \(c_s\ge16L_h\) gives
\[
R_k\le\frac{11}{20}\epsilon+\frac1{160}\epsilon
=\frac{89}{160}\epsilon<\epsilon.
\]
For all sufficiently small \(\epsilon\in(0,1]\), it follows that
\[
T^{\mathrm{KKT}}_\epsilon(\theta(\epsilon),\alpha,0;I)
\le k+1\le K\le C'\epsilon^{-3},\qquad
C':=1+16C_*^2\ell^2c_s\bar\Delta.
\]
Finally, \(\alpha=\kappa\epsilon/(\ell c_s)\), so
\(c_\alpha=\kappa/(\ell c_s)>0\). Our fixed choice of \(\kappa\)
ensures that \(c_\alpha\) depends only on
\(\rho,\gamma,\eta,\tau,\mathbf C\), and \(C'\) depends in addition
on \(\Delta\). Neither constant depends on the instance, the dimensions,
or \(\epsilon\).
\hfill\(\square\)

% Verification status (2026-10-14): the existence and parameter dependence
% of C_* in eq:beta0-gs-rate are taken from GPT Pro's conversion of Li26 v2,
% Section 4.2, (4.19)--(4.23), recorded in
% notes/gpt-proofs/beta0-upper-check-20261013.md, item 7.
% That conversion, including its fixed averaging factor, has not been
% independently verified line by line. Reading the cited statements and
% estimates does not constitute such a verification.
